% Source: https://artofproblemsolving.com/community/c3h1634220_need_help
% Stable identity: aops-topic-1634220-q01
From the first equation $y=7-2x$. Substituting into the second,
\[
7-2x-kx=3\iff(2+k)x=4 .
\]
If $k\ne-2$ this gives the single value $x=\dfrac{4}{2+k}$, and then $y=7-2x$, so there is exactly one solution. If $k=-2$ the equation reads $0=4$ and there is no solution (the two lines $y=7-2x$ and $y=-2x+3$ are parallel). So the system has exactly one solution for $\boxed{\text{every }k\ne-2}$. For example, $k=0$ gives $x=2$, $y=3$.
